% ===== CipherFlow Progress Report - Version 2 =====
\documentclass[a4paper,12pt]{report}

\usepackage[french]{babel}
\usepackage[T1]{fontenc}
\usepackage[utf8]{inputenc}

\usepackage{geometry}
\geometry{left=2.5cm,right=2.5cm,top=2cm,bottom=3.5cm}

\usepackage{graphicx}
\usepackage{float}
\usepackage{hyperref}
\usepackage{wallpaper}

\usepackage{booktabs}
\usepackage[table]{xcolor}
\usepackage{colortbl}

\usepackage{array}
\usepackage{tabularx}
\usepackage{longtable}
\usepackage{pdflscape}
\usepackage{enumitem}
\usepackage{listings}
\usepackage{tcolorbox}
\tcbuselibrary{skins, breakable}

\definecolor{FSTblue}{RGB}{0,102,204}
\definecolor{codebg}{RGB}{245,245,245}
\definecolor{codeframe}{RGB}{200,200,200}

\lstdefinestyle{CodeStyle}{
  basicstyle=\ttfamily\small,
  backgroundcolor=\color{codebg},
  frame=single,
  rulecolor=\color{codeframe},
  breaklines=true,
  showstringspaces=false,
  keywordstyle=\color{FSTblue}\bfseries,
  stringstyle=\color{teal},
  commentstyle=\color{gray}\itshape,
  numbers=left,
  numberstyle=\tiny\color{gray},
  numbersep=5pt,
  xleftmargin=10pt,
}

\frenchbsetup{StandardLists=true}

\newcommand{\question}[1]{\section*{#1}\addcontentsline{toc}{section}{#1}}

\newcommand{\screenshot}[3][0.75]{%
  \IfFileExists{./assets/images/#2}{%
    \begin{figure}[H]\centering
      \includegraphics[width=#1\textwidth]{./assets/images/#2}
      \caption{#3}
    \end{figure}}{%
    \begin{figure}[H]\centering
      \fbox{\parbox[c][5cm][c]{#1\textwidth}{\centering
        \textbf{CAPTURE A INSERER}\\[1ex]\texttt{assets/images/\detokenize{#2}}}}
      \caption{#3}
    \end{figure}}%
}

%%%%%%%%%%%%%%%%%%%%%%%%%%%%
% Custom Title Page Setup  %
%%%%%%%%%%%%%%%%%%%%%%%%%%%%
\newcommand*{\makeCustomTitle}{
    \newgeometry{margin=2.5cm}
    \begin{titlepage}
        \ThisLRCornerWallPaper{1.2}{./assets/branding/background.jpg}
        \setlength{\parindent}{0pt}

        \begin{minipage}[t][0.25\textheight]{\textwidth}
            \begin{center}
                \includegraphics[height=0.15\textheight]{./assets/branding/LOGO_FSTM.png}
            \end{center}
            \vspace{2ex}
            \begin{center}
                {
                    \scshape
                    {\LARGE Universit\'{e} Cadi Ayyad\\}
                    {\Large Facult\'{e} des Sciences et Techniques de Marrakech\\[1ex]}
                    {\large D\'{e}partement d'Informatique}
                }
            \end{center}
        \end{minipage}

        \begin{minipage}{\textwidth}
            \vspace{0.15\textheight}
            \begin{center}
                {\bfseries\huge Etat d'Avancement (Version 2) : CipherFlow}
            \end{center}
            \vspace{3ex}
        \end{minipage}

        \begin{center}
            {\Large Module : Developpement Frontend}

            \vfill

            \begin{tabular}{@{}p{0.49\textwidth}@{\hspace{0.02\textwidth}}p{0.49\textwidth}@{}}
                \begin{minipage}[t]{0.49\textwidth}
                    \large
                    \textbf{\'{E}tudiants :}\\
                    Mohamed \textsc{EL BOUCHIR}\\
                    Bilal \textsc{AKENSOUSS}\\
                    Mohammed \textsc{ALHAYANE}\\
                    Salam \textsc{ELOUZI}
                \end{minipage}
            \end{tabular}

            \vfill

            {\large\today}
        \end{center}
    \end{titlepage}
    \restoregeometry
}

\begin{document}

\makeCustomTitle

\tableofcontents
\newpage

\chapter*{CipherFlow : Tableau de Bord Anime de Transformation de Donnees}
\addcontentsline{toc}{chapter}{CipherFlow : Tableau de Bord Anime de Transformation de Donnees}

CipherFlow est une application web interactive de visualisation d'algorithmes
cryptographiques et de transformation de donnees, inspiree de CyberChef.
L'objectif est d'offrir une plateforme educative ou chaque etape d'un pipeline
est visualisee pas a pas grace a des animations declaratives fluides. Cette deuxieme
version de l'etat d'avancement integre les ameliorations apportees au premier draft,
le detail des taches accomplies ce matin, ainsi que les perspectives planifiees
par l'equipe sous la conduite du chef de projet.

\question{Ameliorations apportees au premier draft et taches accomplies ce matin}

Durant la seance de ce matin, l'equipe a franchi plusieurs etapes cles pour faire
evoluer le projet au-dela du premier draft :

\begin{itemize}
  \item \textbf{Mise en place et amelioration du depot GitHub} :
        Le code source a ete organise, initialise sous Git et heberge sur le depot officiel :\\
        \url{https://github.com/mohamed-el-bouchir/cipherflow}
  \item \textbf{Ajout d'une description officielle du projet} :
        Une description claire et percutante a ete definie sur GitHub (section About)
        pour presenter instantanement la vocation de CipherFlow aux visiteurs.
  \item \textbf{Redaction d'un README detaille} :
        Le fichier \texttt{README.md} a ete cree et structure de maniere exhaustive.
        Il comprend une presentation du concept, les fonctionnalites principales, la liste
        complete de la pile technique, l'arborescence structuree du projet, la documentation
        des 4 algorithmes existants, un guide pour ajouter de nouvelles operations, et
        les instructions pour l'installation et l'execution locale.
  \item \textbf{Retouches et ajustements de l'interface utilisateur (UI)} :
        L'interface a ete amelioree au niveau de l'agencement visuel : alignement plus
        net des panneaux de saisie et de sortie, lisibilite accrue du texte en cours
        de transformation et ajustement des espacements de la scene d'animation.
  \item \textbf{Stabilisation de la configuration technique} :
        Nettoyage et fiabilisation des dependances du projet dans le fichier \texttt{package.json}
        pour garantir une installation et un demarrage fluides sans conflits de versions.
\end{itemize}

\screenshot[0.45]{github-repository.png}{Apercu du depot GitHub officiel avec documentation et README}

\question{Presentation du projet et environnement technique}

CipherFlow permet aux utilisateurs de composer une chaine d'operations (chiffrement,
encodage, transformations binaires et manipulation de texte) tout en observant
le flux interne de calcul via un lecteur pas a pas.

La pile technique utilisee comprend :
\begin{itemize}
  \item \textbf{React 18 et TypeScript} : composants modulaires et typage statique rigoureux.
  \item \textbf{Bootstrap} : mise en page structuree et responsive design de l'interface.
  \item \textbf{Motion (anciennement Framer Motion)} : moteur d'animation declarative.
  \item \textbf{Vite} : serveur de developpement et assemblage optimise des ressources.
  \item \textbf{Lucide React} : bibliotheque d'icones vectorielles.
\end{itemize}

\screenshot[0.9]{workbench-initial.png}{Interface de CipherFlow avec les ameliorations d'agencement}

\question{Architecture du workbench et pipeline d'operations}

L'interface de travail s'organise selon cinq zones principales :
\begin{enumerate}
  \item \textbf{TopBar} : barre superieure permettant de basculer entre l'atelier
        (Workbench) et les recettes predefinies (Presets), avec controle global de lecture.
  \item \textbf{InputPanel} : zone de saisie du texte initial a transformer.
  \item \textbf{RecipePipeline} : pipeline de composition permettant d'ajouter, d'activer,
        de desactiver, de deplacer ou de supprimer des operations, avec prise en charge
        du mode dechiffrement et de la sauvegarde de recettes.
  \item \textbf{AnimationStage} : scene centrale affichant le composant \texttt{Visualizer}
        propre a chaque operation, dotee d'un lecteur pas a pas et d'un selecteur de vitesse.
  \item \textbf{OutputPanel} : affichage de la sortie intermediaire de l'etape active
        et du resultat final global.
\end{enumerate}

\screenshot[0.9]{pipeline-caesar.png}{Animation pas a pas de l'operation Caesar Cipher}

\screenshot[0.9]{pipeline-multistep.png}{Vue du pipeline multi-etapes avec calcul des sorties}

\question{Systeme d'extension des operations}

Le coeur de CipherFlow repose sur un registre modulaire (\texttt{src/algorithms/registry.ts}).
Chaque operation implemente une interface standard permettant d'ajouter de nouvelles
transformations sans impacter le reste du systeme :

\begin{lstlisting}[style=CodeStyle, language={[Sharp]C}, caption={Interface AlgorithmDefinition (TypeScript)}]
export interface AlgorithmDefinition<TState = any> {
  id: string;
  name: string;
  category: AlgorithmCategory;
  summary: string;
  animationPattern: string;
  complexity: string;
  paramSchema: ParamField[];
  defaultParams: Record<string, any>;
  computeTrace: (input: string, params: Record<string, any>)
    => AlgorithmTrace<TState>;
  invertParams?: (params: Record<string, any>) => Record<string, any>;
  Visualizer: ComponentType<VisualizerProps<TState>>;
}
\end{lstlisting}

Quatre operations initiales servent de base de reference :
\begin{itemize}
  \item \textbf{Caesar Cipher} (\texttt{caesar.tsx}) : decalage de lettres dans l'alphabet.
  \item \textbf{Base64 Encode/Decode} (\texttt{base64.tsx}) : conversion en blocs Base64.
  \item \textbf{XOR Cipher} (\texttt{xor.tsx}) : operation binaire octet par octet avec cle.
  \item \textbf{Reverse Text} (\texttt{reverse.tsx}) : inversion et permutation de casse.
\end{itemize}

\question{Direction du projet et prochaines etapes}

Le projet est dirige par \textbf{Mohamed El Bouchir} en qualite de chef de projet,
qui assure la coordination generale des travaux avec l'ensemble des membres de l'equipe :
Bilal Akensouss, Mohammed Alhayane et Salam Elouzi.

A partir de cette version socle, l'equipe va realiser les evolutions prioritaires suivantes :

\begin{enumerate}
  \item \textbf{Creation d'une Landing Page (Page d'accueil)} :\\
        Concevoir une page d'accueil moderne pour introduire le projet, expliciter
        les concepts cles de CipherFlow et proposer un point d'entree clair vers l'atelier.

  \item \textbf{Integration d'un popup de premier usage (First-use Tutorial Popup)} :\\
        Mettre en place une fenetre explicative s'affichant lors de la premiere visite
        pour guider l'utilisateur a travers le pipeline, les parametres d'operations et
        les fonctionnalites de la scene d'animation.

  \item \textbf{Amelioration continue de l'interface et de l'experience utilisateur (UI/UX)} :\\
        Poursuivre l'harmonisation visuelle avec Bootstrap, enrichir les transitions
        animees et fluidifier les manipulations de composition.

  \item \textbf{Ajout de nouvelles operations de transformation} :\\
        Etendre le catalogue d'operations avec de nouveaux algorithmes de chiffrement,
        d'encodage et de traitement de texte, accompagnes de leurs animations dediees.
\end{enumerate}

\bigskip

Cette deuxieme version atteste des progres concrets accomplis ce matin : le projet
est proprement documente sur GitHub, l'interface est amelioree, et la direction des
prochaines etapes est clairement definie.

\end{document}
